\section{Introduction}
\label{sec:intro}

Access control for the Internet of Things (IoT) is no longer a database
permission problem. A modern deployment authenticates battery-powered sensors,
issues commands to valves, locks and motors, and continuously re-evaluates who
may do what under changing load, link quality and physical state
\cite{ravidas2019access,qiu2020survey}. Context-adaptive access control (CAAC)
responds by making the security posture itself a function of context: a quiet
channel with a healthy battery may use a lighter handshake, while congestion or
anomalous telemetry triggers stronger authentication and deeper inspection
\cite{atlam2017overview,rath2017adaptive,genc2019context,arfaoui2019context}.
The promise is attractive because the alternative---a single high posture---is
expensive in energy, latency and false denials on constrained devices
\cite{elgendi2019protecting,elhajj2019survey}.

The same adaptability is an attack surface. An adversary who can influence the
context channel can make an unsafe posture look rational: spoof a low battery,
suppress anomaly evidence, or replay a benign-looking neighbourhood to drive the
policy below the level the physical process requires. Sensor spoofing and false
data injection are well documented in cyber-physical systems
\cite{barua2020special,li2019optimal,ye2020summation}, and denial-of-sleep
attacks show that forcing a device to over-react is equally effective
\cite{brownfield2005wireless,raymond2009effects}. Byzantine nodes can bias an
aggregate without controlling a majority \cite{lamport1982byzantine}, and a
periodic attacker can make an adaptive controller oscillate between postures
until it exhausts energy or its own hysteresis \cite{kephart2003vision,cheng2009software}.
Existing CAAC designs optimise expected utility against a point estimate of
risk and treat security strength as one more soft cost
\cite{lee2017situational,moura2017sotraace,hussainalnaji2022cab}. Under a
coherent spoof, the point estimate moves in the wrong direction, and because
nothing in the objective is non-negotiable, the optimiser follows it.

This paper starts from a different decomposition. \emph{Adaptation may choose
among safe actions, but it may not choose whether an action is safe.} We define
an admissible set $\Asafe(\op,\attack_t)$ from unconditional invariants---the
actuation envelope I1, the safety-traffic envelope I2, a belief-conditional
restriction I3, and per-operation failure semantics I4---and optimise mission
utility only inside that set. The invariants do not depend on the threat belief,
so a compromised or merely unlucky estimator cannot argue its way below them.
When the set would be empty, the enforcement point executes a precomputed safe
action: fail-closed for confidentiality and integrity operations, and
fail-safe-allow (authenticated emergency function only) where denial would
itself violate a physical safety requirement. This is runtime assurance in the
sense of Simplex and shield synthesis \cite{sha2001simplicity,wu2017safety,wu2019shield,mehmood2022black,cho2025timed},
but applied to an access-control action lattice rather than to a single
controller switch.

Inside the admissible set we keep uncertainty explicit. A forward hidden Markov
model maintains a belief over five latent threat states, and four context
modalities are fused with weights that fall when a modality's posterior drifts
from the device-telemetry anchor. The monitor therefore distinguishes ``the
plant is calm'' from ``the telemetry is lying'' as far as the evidence allows,
and the belief is scored with strictly proper rules rather than accuracy
\cite{guo2017calibration,gneiting2007strictly}. The controller adds asymmetric
escalation, a dwell band and a utility margin, so a resonant attacker cannot
convert belief noise into posture thrashing. Cryptographic enforcement uses
AEAD and lightweight authenticated encryption at the action level
\cite{dobraunig2021ascon,cagua2025implementation}; the contribution is the
decision layer that selects among those mechanisms, not a new cipher.

We evaluate the design in two complementary ways. The primary evaluation is a
calibrated discrete-time simulation of a 24-node deployment with 2{,}000 epochs,
10 paired seeds, five latent states, four context modalities, eight scenarios
(seven attacks plus benign) and ten policies. It is a parameterised simulation,
not a replay of a public dataset: every scenario is generated from an explicit
emission and transition model, and the same stream is replayed to every policy
so that comparisons are paired. The network-intensity prior is calibrated
against the public IoT-23 capture \cite{garcia2020iot}, and we report the
measured benign-to-malicious contrast and its per-capture spread rather than
claiming raw-trace training. Hardware latency, energy and footprint are analytical
estimates from datasheet-level operation costs; they are labelled as modelled
throughout and are not silicon measurements.

\noindent\textbf{Contributions.}
\begin{itemize}
  \item \textbf{C1 (RQ1): a constrained CAAC formulation with hard invariants.}
  We define the admissible set, the per-operation failure semantics, and a
  constrained objective in which invariants are not tradeable. Proposition~1
  shows that the invariant filter is dominant and that a safe action always
  exists; Proposition~2 shows that a granted actuation or safety request can never
  fall below the structural floor, for any belief, including an adversarially
  chosen one
  (\cref{sec:design}).
  \item \textbf{C2 (RQ3): continuous risk belief with cross-modal consistency
  weights.} A forward HMM is fused with trust weights derived from
  modality-posterior disagreement and a cross-node dispersion signal. The belief
  is well calibrated under DoS, downgrade, energy-inflation, Byzantine and
  adaptive attacks, because the attested device anchor and the cross-node
  dispersion survive those poisons. When the anchor itself is subverted---an
  explicitly out-of-model \texttt{tee} stress test---the belief collapses, which
  is why the safety guarantee is carried by I1/I2 rather than by detection
  (\cref{sec:design,sec:evaluation}).
  \item \textbf{C3 (RQ5): fail-safe enforcement with bounded adaptation.}
  Asymmetric escalation, a dwell band and a utility margin reduce policy
  switches by more than half relative to the same controller without hysteresis,
  with essentially unchanged utility and lower residual exposure
  (\cref{sec:design,sec:evaluation}).
  \item \textbf{C4 (RQ1--RQ4): decision-level evaluation under adaptive
  poisoning.} We report invariant violations, residual malicious grants,
  false denials, mission utility, calibration and safety-request denial across
  seven attack scenarios and ten policies, with paired non-parametric tests and
  Holm correction (\cref{sec:evaluation}).
  \item \textbf{C5 (RQ5): modelled systems cost on three tiers.} We give an
  analytical cost model for Class-1 MCU, Class-2 MCU and gateway tiers and report
  p50/p99 decision latency, energy per decision and footprint. The full
  controller is cheaper than a static-high posture in both latency and energy
  because it selects the cheapest admissible action most of the time
  (\cref{sec:implementation,sec:evaluation}).
  \item \textbf{C6 (RQ1, RQ5): trace-derived calibration and floor sensitivity.}
  We calibrate the simulator's network-intensity prior against IoT-23, report the
  fitted mean-log-rate contrast with a per-capture range, and measure how the
  strength of the actuation floor trades residual exposure against utility and
  energy
  (\cref{sec:evaluation}).
\end{itemize}

The rest of the paper is organised as follows. \Cref{sec:background} introduces
the access-control and self-adaptive background. \Cref{sec:threat-model} states
the system and threat model, \cref{sec:design} gives the formal design and
algorithms, and \cref{sec:implementation} describes the simulator and cost model.
\Cref{sec:evaluation} answers RQ1--RQ5, \cref{sec:discussion} discusses
limitations and threats to validity, \cref{sec:related} relates the work to prior
art, and \cref{sec:conclusion} concludes.
